% from ega4-2-fr.tex
\oldpage[IV-2]{173}
Comme l'intersection des $\mathfrak r'_j$ est réduite à 0, on en déduit que
$B'_{\mathfrak p'}$ lui-même est noethérien, en raison du lemme classique
suivant:

\begin{lemma}[6.14.1.4]
\label{IV.2.s6.14.1.4-fr}
Soient $R$ un anneau, $\mathfrak a$ et $\mathfrak b$ deux idéaux de $R$; si
$R/\mathfrak a$ et $R/\mathfrak b$ sont noethériens, il en est de même de
$R/(\mathfrak a\cap\mathfrak b)$.
\end{lemma}

\begin{proof}
En effet, soit $\mathfrak c$ un idéal de $R$ tel que
$\mathfrak a\cap\mathfrak b\subset\mathfrak c$; on a donc
$\mathfrak c\cap\mathfrak b\subset\mathfrak a\cap\mathfrak c\subset
\mathfrak c$; or $\mathfrak c/(\mathfrak a\cap\mathfrak c)$ est un
$R$-module isomorphe à $(\mathfrak a+\mathfrak c)/\mathfrak a$, donc à un
idéal de $R/\mathfrak a$, et par suite est de type fini; d'autre part
$(\mathfrak a\cap\mathfrak c)/(\mathfrak a\cap\mathfrak b)$ est un
sous-$R$-module de $\mathfrak a/(\mathfrak a\cap\mathfrak b)$, et ce
dernier est isomorphe à $(\mathfrak a+\mathfrak b)/\mathfrak b$, donc est
de type fini comme idéal de $R/\mathfrak b$; donc
$(\mathfrak a\cap\mathfrak c)/(\mathfrak a\cap\mathfrak b)$ est aussi un
$R$-module de type fini, et il en est finalement de même de
$\mathfrak c/(\mathfrak a\cap\mathfrak b)$.
\end{proof}

Notons d'autre part que $\operatorname{Spec}(A)$ se compose du point fermé
$\mathfrak m$ et du point générique $(0)$, de corps résiduels $k$ et $K$
respectivement; comme dans le cas où nous nous sommes placés, le corps des
fractions de $B$ et son corps résiduel sont respectivement isomorphes à $K$
et $k$, les fibres du morphisme
$u:\operatorname{Spec}(B'_{\mathfrak p'})\to\operatorname{Spec}(B)$ au
point fermé et au point générique de $\operatorname{Spec}(B)$ sont
respectivement isomorphes aux fibres du morphisme
$\operatorname{Spec}(A'_{\mathfrak p'})\to\operatorname{Spec}(A)$ au point
fermé et au point générique de $\operatorname{Spec}(A)$ (I, 3.6.4), donc
sont géométriquement normales par hypothèse; comme d'autre part le
morphisme $u$ est plat, on en conclut qu'il est normal (6.8.1). Mais comme
$B$ et $B'_{\mathfrak p'}$ sont noethériens et que $B$ est normal, on
déduit de (6.5.4) que $B'_{\mathfrak p'}$ est normal.

Cela étant, $B$ est réunion de la famille filtrante croissante de ses
sous-$A$-algèbres finies $A_\alpha$; par platitude, $B'$ est réunion de la
famille filtrante croissante des
$A'_\alpha=A_\alpha\otimes_A A'$; si l'on pose
$\mathfrak p'_\alpha=\mathfrak q'\cap A'_\alpha$,
$B'_{\mathfrak q'}$ est aussi réunion de la famille filtrante croissante
des $(A'_\alpha)_{\mathfrak p'_\alpha}$ (5.13.3). Désignons par $L''$ le
composé direct des corps $L'_j$ tels que
$\mathfrak q'_j\not\subset\mathfrak p'$; alors pour tout $\alpha$,
$(A'_\alpha)_{\mathfrak p'_\alpha}$ est contenu dans
$K'_{\mathfrak p'}$, et l'anneau
$\overline{B'_{\mathfrak q'}}=L''\times B'_{\mathfrak q'}$ est donc réunion
des anneaux
$L''\times(A'_\alpha)_{\mathfrak p'_\alpha}=
(\overline{A'_\alpha})_{\mathfrak p'_\alpha}$. Mais chacun des $A_\alpha$
est local, noethérien, intègre et de dimension 1, et le morphisme
$\operatorname{Spec}(A'_\alpha)\to\operatorname{Spec}(A_\alpha)$ est
normal (6.8.2), donc on peut lui appliquer le lemme (6.14.1.2) et l'on a
\[
A'_\alpha=K'\cap(\overline{A'_\alpha})_{\mathfrak p'_\alpha}
\]
pour tout $\alpha$; prenant la limite inductive de chacun des deux membres,
il vient
\[
B'=K'\cap\overline{B'_{\mathfrak q'}}.
\]
Mais $\overline{B'_{\mathfrak q'}}$, composé direct des anneaux normaux
$L''$ et $B'_{\mathfrak q'}$, est normal, et comme il en est de même de
$K'$, le lemme (6.14.1.1) montre que $B'$ est normal. C.Q.F.D.
\end{proof}

\begin{corollary}[6.14.2]
\label{IV.2.s6.14.2-fr}
Soient $k$ un corps, $X$ un $k$-préschéma normal (non nécessairement
localement noethérien). Alors, pour toute extension séparable $k'$ de $k$,
$X_{(k')}=X\otimes_k k'$ est normal.
\end{corollary}

\begin{proof}
On sait en effet (6.7.6) que le morphisme
$\operatorname{Spec}(k')\to\operatorname{Spec}(k)$ est alors normal.
\end{proof}

\begin{corollary}[6.14.3]
\label{IV.2.s6.14.3-fr}
Soient $k$ un corps; $X$, $Y$ deux $k$-préschémas intègres et normaux, dont
les corps de fonctions rationnelles $K=R(X)$, $L=R(Y)$ sont des extensions
séparables de $k$. Alors le préschéma $X\times_k Y$ est normal.
\end{corollary}

\begin{proof}
